\documentclass[10pt,a4paper]{amsart}
\usepackage[margin=19mm]{geometry}
\usepackage{amsmath,amssymb,mathtools}
\usepackage[scr=rsfs,cal=euler]{mathalfa}
\usepackage{xfrac}
\usepackage[unicode,hidelinks,bookmarks=false]{hyperref}
\newcommand{\ie}{i.e.\ }
\newcommand{\cf}{cf.\ }
\newcommand{\resp}{respectively\ }
\newcommand{\Subsection}{\subsection}
\providecommand{\nobreakdash}{\nobreak\hbox{-}}
\setlength{\emergencystretch}{2em}
\multlinegap0pt
\usepackage{amssymb}
\usepackage{mathtools}
\newtheorem{theorem}{Theorem}
\newcommand{\Lr}{\operatorname{L}}
\title{On a conjecture due to Kanade related to Nahm sums}
\author{Cetin Hakimoglu-Brown}
\date{Source: arXiv:2603.00836v3 (CC BY 4.0)}
\begin{document}
\maketitle

\noindent\emph{Source and licence.} On a conjecture due to Kanade related to Nahm sums. Version arXiv:2603.00836v3 (CC BY 4.0), distributed by arXiv under the Creative Commons Attribution 4.0 International licence, http://creativecommons.org/licenses/by/4.0/, as recorded in the arXiv licence metadata for this exact version and shown on the abstract page https://arxiv.org/abs/2603.00836v3. Full licence notice: \texttt{licenses/arxiv-2603.00836-CC-BY-4.0.txt}.\par\smallskip

\noindent\textit{Editorial scope.} Counted unit: the article's theorem thm:kanade (source0.tex lines 205--207). The article's complete original proof of it is delivered with it (source0.tex lines 209--241, one \texttt{proof} environment of 33 lines). No mathematics is written, repaired or re-derived: the statement and the proof are the article's own lines, in the article's own notation, and no retained line is substituted or re-flowed.  Retained uncounted as the source-local context the proof needs: lines 100--102; lines 187--190; lines 193--196; lines 202--204.  Nothing was omitted from any retained span, so the package carries no omission record.  No retained line contains a citation, so the wrapper carries no bibliography and no bibliography entry.  No retained line contains a figure environment, an image-inclusion command or drawing code, so no image is delivered and the package declares no asset dependency.  Editorial formatting only, in addition: an AMS \texttt{amsart} preamble that restates the article's own declarations and macros used by the retained lines, namely the environments \texttt{theorem} on separate counters, together with the standard packages those lines need.  The retained environments print in this document as Theorem 1 in order of appearance, whereas the article prints the same items with the numbering of its own full document.  The complete native source of the article as distributed in the arXiv e-print for this exact version is bundled unchanged as \texttt{sources/195518/source0.tex}.
\begin{equation}\label{eq:kanade}
\Lr(Q_1)+\frac13\Lr(Q_2^3)=\frac{4\pi^2}{27}.
\end{equation}

\begin{equation}\label{eq:reduced1}
\Lr(x^3)-3\Lr(x^2)-\Lr(x)
=\Lr(z^2)-\Lr(y^2)-\Lr(1).
\end{equation}

\begin{equation}\label{eq:reduced2}
2\Lr(y)-\frac12\Lr(y^2)
=\Lr(z)-\frac12\Lr(z^2)+\Lr(x)+\Lr(b).
\end{equation}

\begin{equation}\label{eq:loxton}
\Lr(x^3)-3\Lr(x^2)-3\Lr(x)=-\frac73\Lr(1).
\end{equation}

\begin{theorem}\label{thm:kanade}
Kanade's identity~\eqref{eq:kanade} holds.
\end{theorem}

\begin{proof}
The $x^3$ term can now be removed.  Subtracting
\eqref{eq:loxton} from~\eqref{eq:reduced1} and rearranging gives:
\begin{equation}\label{eq:diffsq}
\Lr(y^2)-\Lr(z^2)=\frac43\Lr(1)-2\Lr(x).
\end{equation}
Rearranging~\eqref{eq:reduced2} and then using~\eqref{eq:diffsq} gives:
\begin{align}
2\Lr(y)-\Lr(b)
&=\Lr(x)+\Lr(z)+\frac12\bigl(\Lr(y^2)-\Lr(z^2)\bigr)\notag\\
&=\Lr(z)+\frac23\Lr(1).\label{eq:intermediate}
\end{align}
Since $y+z=1$, Euler's relation gives $\Lr(z)=\Lr(1)-\Lr(y)$, and hence:
\begin{equation}\label{eq:keykanade}
3\Lr(y)-\Lr(b)=\frac53\Lr(1).
\end{equation}
This is the entire dilogarithmic cancellation.  Applying Euler's relation
once more:
\begin{equation}\label{eq:kanadefinish}
\Lr(y)+\frac13\Lr(1-b)
 =\frac13\bigl(3\Lr(y)-\Lr(b)+\Lr(1)\bigr)
 =\frac89\Lr(1)=\frac{4\pi^2}{27}.
\end{equation}

It remains only to identify the arguments with Kanade's notation.  From
the definitions:
\[
Q_1=\frac{1}{1+x}=y,
\qquad
Q_2^3=1-\frac{x^2}{1+x}=1-b.
\]
Thus \eqref{eq:kanadefinish} is exactly \eqref{eq:kanade}.
\end{proof}

\end{document}
